\documentclass[11pt]{article}
\usepackage[T1]{fontenc}
\usepackage{lmodern,amsmath,amssymb,amsthm,booktabs,array}
\usepackage[margin=28mm]{geometry}
\usepackage[colorlinks=true,linkcolor=blue,urlcolor=blue,citecolor=blue]{hyperref}
\newtheorem{theorem}{Theorem}
\newtheorem{corollary}[theorem]{Corollary}
\newcommand{\NN}{\mathbb N}
\newcommand{\floor}[1]{\left\lfloor #1\right\rfloor}
\title{Resolving the Hermite coefficient conjectures in\\
OEIS A381423, Comments (revision 8)}
\author{Haobo Ma\\\small ChronoAI Pte Ltd, Singapore
\and Wenlin Zhang\\\small National University of Singapore\\
\small ChronoAI Pte Ltd, Singapore}
\date{3 November 2026}
\begin{document}
\maketitle
\begin{abstract}
We prove Mike Sheppard's conjecture in OEIS A381423, Comments, revision 8,
that the exponent of the largest signed coefficient of the physicists'
Hermite polynomial has successive differences in $\{1,-3\}$.
A coefficient-ratio calculation gives an explicit formula and proves
uniqueness, square-root growth, and the two asserted difference-block lengths.
We give exact quadratic indices for the longer blocks.
The corresponding absolute-coefficient calculation proves the conjecture in
A381524 and the three formulas for equal absolute maxima in A381737.
It also gives the absolute mode explicitly, locates its two difference-block
types, and determines the asymptotic ratios of successive maximal coefficients.
\end{abstract}

\section{The signed maximum}
OEIS A381423, Comments, revision 8, conjectures
``Differences are either 1 or -3''~\cite{signed}.
Its other observations concern uniqueness, blocks of lengths four and five,
their spacing, and growth. The ratio calculation below proves uniqueness,
the block lengths, and the growth assertion. It specifies spacing by exact
starting indices and intervening block counts.
The sign convention matters: absolute maxima form a different sequence.
The coefficient array is A060821~\cite{triangle}; removing parity zeros gives
A059343~\cite{nonzero}. We use the physicists'
normalization, whose finite expansion is DLMF 18.5.13~\cite{dlmf}:
\begin{equation}\label{eq:hermite}
 H_n(x)=n!\sum_{k=0}^{\floor{n/2}}
       \frac{(-1)^k(2x)^{n-2k}}{k!(n-2k)!},\qquad n\in\NN.
\end{equation}
Write $a(n)$ for the least $j\in\{0,\ldots,n\}$ at which
$[x^j]H_n(x)$ is maximal, comparing signed coefficients.
The least-index convention agrees with the programs in the posing entry;
the theorem below shows that the maximum is in fact unique.
The two-valued difference assertion also has a formal proof in the
\href{https://github.com/the-omega-institute/trureturing/blob/a2d277ea0980f1230eff2c7e27838372effccb30/D5/S3/Analytic/Polynomial/HermiteSignedMode.lean}{pinned Lean source}.
That proof dates from 2 November 2026; the formula and corollaries below
give a direct coefficient-ratio treatment of the neighboring assertions.

\begin{theorem}\label{thm:signed}
For every $n\geq0$, the largest signed coefficient of $H_n$ is unique, and
\begin{equation}\label{eq:closed}
 q(n)=\floor{\frac{2n+7-\sqrt{8n+21}}8},\qquad a(n)=n-4q(n).
\end{equation}
Equivalently, $a(n)$ is the least nonnegative integer $j\equiv n\pmod4$
such that
\begin{equation}\label{eq:F}
 F(j):=j^2+7j+7\ \geq\ 2n.
\end{equation}
Moreover, $a(n+1)-a(n)\in\{1,-3\}$ for every $n\geq0$.
\end{theorem}

\begin{proof}
By~\eqref{eq:hermite}, the positive coefficients are exactly
\[
 C_{n,r}=\frac{n!2^{n-4r}}{(2r)!(n-4r)!},
 \qquad 0\leq r\leq\floor{n/4}.
\]
All other coefficients are negative or zero. Thus the signed maximum
is among these positive numbers, including $C_{n,0}=2^n$.
For $1\leq r\leq\floor{n/4}$ put $j=n-4r$. Cancelling positive
factorials gives
\begin{equation}\label{eq:ratio}
 \frac{C_{n,r}}{C_{n,r-1}}
 =\frac{(j+1)(j+2)(j+3)(j+4)}{16(2r)(2r-1)}.
\end{equation}
This ratio strictly decreases with $r$: its numerator strictly decreases
and its positive denominator strictly increases.
It never equals one. Indeed, put $X=j^2+5j+5\geq5$ and $Y=8r-2\geq6$.
The numerator minus denominator in~\eqref{eq:ratio} is
\begin{equation}\label{eq:squares}
 X^2-Y^2+3.
\end{equation}
If $X\geq Y$, this is positive. If $X<Y$, integrality gives
$X^2-Y^2+3\leq X^2-(X+1)^2+3=2-2X<0$.
Consequently
\[
 C_{n,r}>C_{n,r-1}
 \ \Longleftrightarrow\ j^2+5j+7\geq8r
 \ \Longleftrightarrow\ F(j)\geq2n.
\]
As $F$ is strictly increasing on the nonnegative integers, the coefficients
increase strictly through the last qualifying $r$ and decrease strictly
thereafter. For $r=0$, the inequality $F(n)\geq2n$ always holds.
There is therefore a unique maximum, at the largest qualifying $r=q(n)$.
This proves~\eqref{eq:F}, including the restriction $j\leq n$ because
$j=n$ qualifies.

For $j\geq0$, condition~\eqref{eq:F} is equivalent to
\[
 j\geq\alpha_n:=\frac{\sqrt{8n+21}-7}{2}.
\]
For $n\geq4$, $\alpha_n>0$, so the largest allowable $r$ is
$\floor{(n-\alpha_n)/4}$. This is~\eqref{eq:closed}.
For $0\leq n\leq3$, both the coefficient definition and that formula give
$q(n)=0$ and $a(n)=n$.

It remains to compare adjacent degrees. If $r$ qualifies at degree $n$
and $j=n-4r$, then it qualifies at degree $n+1$, since
\[
 F(j+1)-2(n+1)=F(j)-2n+2j+6>0.
\]
Thus $q(n+1)\geq q(n)$. Conversely, if $r+1$ qualifies at degree $n+1$
and $j=n+1-4(r+1)\geq0$, then $r$ qualifies at degree $n$, since
\[
 F(j+3)-2n=F(j)-2(n+1)+6j+32>0.
\]
Taking $r=q(n+1)-1$ when $q(n+1)>0$ shows that
$q(n+1)\leq q(n)+1$; the case $q(n+1)=0$ is immediate.
Hence $a(n+1)-a(n)=1-4(q(n+1)-q(n))$ is $1$ or $-3$.
\end{proof}

The largest signed coefficient, indexed separately in A277280~\cite{signedvalue},
is therefore given exactly by
\begin{equation}\label{eq:maxvalue}
 B(n)=\frac{n!2^{a(n)}}{(2q(n))!a(n)!}.
\end{equation}
For example, $H_5(x)=32x^5-160x^3+120x$ gives $q(5)=1$,
$a(5)=1$, and $B(5)=120$.

\section{Blocks, spacing, and growth}
Set $n_0=0$. For $r\geq1$ define
\begin{equation}\label{eq:thresholds}
 j_r=\min\{j\geq0:j^2+5j+7\geq8r\},\qquad n_r=4r+j_r.
\end{equation}
The threshold characterization in the proof shows that $q(n)\geq r$
if and only if $n\geq n_r$. Thus $n_r$ is the first degree with mode index $r$.

\begin{corollary}\label{cor:blocks}
The difference sequence is a concatenation of
$(1,1,1,-3)$ and $(1,1,1,1,-3)$. The first block has length five.
For $j\geq1$, let
\begin{equation}\label{eq:R}
 R_j=\floor{\frac{j^2+5j+7}{8}}.
\end{equation}
Apart from the first block, the blocks of length five have precisely
the block indices $r=R_j$, where the block indexed by $r$ starts at
difference index $n_r$. Their starting indices are $4R_j+j$.
Between the length-five blocks indexed by $R_j$ and $R_{j+1}$ there are
exactly $R_{j+1}-R_j-1=j/4+O(1)$ length-four blocks.
Furthermore,
\begin{equation}\label{eq:asymptotic}
 \alpha_n\leq a(n)<\alpha_n+4,\qquad
 a(n)=\sqrt{2n}+O(1),\qquad
 \lim_{n\to\infty}\frac{a(n)}{\sqrt{2n}}=1.
\end{equation}
\end{corollary}

\begin{proof}
Put $G(j)=j^2+5j+7$. Since $G(0)=7$, one has $j_1=1$ and $n_1=5$.
The integers $j_r$ are nondecreasing and at least one. Also
$G(j+1)-G(j)=2j+6\geq8$ for $j\geq1$; hence
$j_{r+1}\leq j_r+1$. Therefore $n_{r+1}-n_r\in\{4,5\}$ for $r\geq1$.
For $n_r\leq n<n_{r+1}$, the value of $q(n)$ is $r$.
The corresponding differences are $n_{r+1}-n_r-1$ copies of $1$
followed by $-3$, proving the block description.

The inequality $G(j-1)<8r\leq G(j)$ says exactly that $j_r=j$.
Since $G$ is odd-valued and $G(j)-G(j-1)\geq8$ for $j\geq2$,
the last such $r$ is $R_j$, and the next index has $j_{r+1}=j+1$.
This also holds for $j=1$, directly from $R_1=1$ and $j_2=2$.
Thus the length-five block indices are exactly $R_j$.
The count of intervening blocks follows by subtraction.
Since $R_j=j^2/8+5j/8+O(1)$, their indices and starting degrees
are quadratic in $j$, while the intervening counts grow linearly.
These exact indices specify the spacing without an interval convention.
In particular, the length-four blocks do not have a single constant spacing:
the spacings of \emph{all} block starts are four or five, whereas the
length-five starts have quadratic growth. This is the precise interpretation
of the source's informal linear/quadratic spacing description.

Finally, $a(n)$ is the first integer in its residue class above $\alpha_n$,
so it is less than $\alpha_n+4$. This remains true for $n=0,1,2,3$
by direct substitution. The asymptotic statements follow from
$\alpha_n=\sqrt{2n+21/4}-7/2$.
\end{proof}

\section{The absolute-coefficient siblings}
For completeness, the same method settles the conjecture in A381524,
Comments~\cite{absolute}, and the three conjectures in A381737,
Formula~\cite{ties}. Let $b(n)$ be the least exponent attaining the
largest absolute coefficient of $H_n$.

\begin{theorem}\label{thm:absolute}
The integer $b(n)$ is the least nonnegative $j\equiv n\pmod2$ satisfying
\begin{equation}\label{eq:absolute}
 j^2+5j+2\geq2n.
\end{equation}
For every $n\geq0$, $b(n+1)-b(n)\in\{1,-1\}$.
There are two absolute maxima exactly at the degrees
\begin{equation}\label{eq:ties}
 N_m=2m(m+2)+(m+1)(-1)^{m+1},\qquad m\geq1,
\end{equation}
in increasing order. At every other degree there is one absolute maximum.
\end{theorem}

\begin{proof}
The nonzero absolute coefficients are
$U_{n,k}=n!2^{n-2k}/(k!(n-2k)!)$, for $0\leq k\leq\floor{n/2}$.
For $k\geq1$, with $j=n-2k$, their consecutive ratio is
\[
 \frac{U_{n,k}}{U_{n,k-1}}=\frac{(j+1)(j+2)}{4k}.
\]
These ratios strictly decrease, and a ratio is at least one exactly
when $S(j):=j^2+5j+2\geq2n$. As before, take the largest qualifying
$k=t(n)$, including $k=0$. This proves~\eqref{eq:absolute} and
$b(n)=n-2t(n)$, choosing the smaller exponent at a tie.
Qualification persists from $n$ to $n+1$ because
$S(j+1)-2(n+1)=S(j)-2n+2j+4>0$.
If $k+1$ qualifies at $n+1$ with exponent $j=n+1-2(k+1)$,
then $k$ qualifies at $n$ because
$S(j+1)-2n=S(j)-2(n+1)+2j+8>0$.
Thus $t(n+1)-t(n)\in\{0,1\}$, proving the difference assertion.

Two maxima occur exactly when one ratio equals one; strict decrease
excludes three maxima. Equality gives
$n=(j^2+5j+2)/2$, with $j\equiv n\pmod2$.
The parity condition is equivalent to $4\mid(j+1)(j+2)$,
or $j\equiv2,3\pmod4$. The associated $k=(n-j)/2$
is positive, so every such $j$ supplies a tie.
Putting $j=4s+2$ or $j=4s+3$, for $s\geq0$, yields respectively
$n=8s^2+18s+8$ and $n=8s^2+22s+13$.
The first precedes the second by $4s+5$, and the next first exceeds
the second by $12s+21$. Interleaving them gives~\eqref{eq:ties}.
\end{proof}

The other two formulas in A381737 follow directly: with indexing $m\geq1$,
\begin{align}
 \sum_{m\geq1}N_mx^m
 &=\frac{6x-2x^2}{(1-x)^3}+\frac{x(2+x)}{(1+x)^2}\notag\\
 &=\frac{8x+5x^2+5x^3-x^4-x^5}{(1-x)^3(1+x)^2},\label{eq:gf}\\
 N_m&=N_{m-1}+2N_{m-2}-2N_{m-3}-N_{m-4}+N_{m-5}
 \quad(m\geq6).\label{eq:recurrence}
\end{align}
For the first equality, sum $2m^2+4m$ and $(m+1)(-1)^{m+1}$
separately using the geometric series and its first two derivatives.
Expanding $(1-x)^3(1+x)^2=1-x-2x^2+2x^3+x^4-x^5$
then gives the recurrence by comparing coefficients above degree five.
Its initial values are $8,13,34,43,76$.
The maximal absolute value, A277281~\cite{absolutevalue}, is
\begin{equation}\label{eq:absolutevalue}
 M(n)=\frac{n!2^{b(n)}}{((n-b(n))/2)!b(n)!}.
\end{equation}
For example, at degree eight the equal absolute maxima occur at exponents
two and four, with opposite signs. The signed maximum is unique at exponent four.

\begin{corollary}\label{cor:absolute-blocks}
For every $n\geq0$,
\[
 t(n)=\floor{\frac{2n+5-\sqrt{8n+17}}4},\qquad b(n)=n-2t(n),
 \qquad b(n)=\sqrt{2n}+O(1).
\]
The absolute-mode difference sequence consists of $(1,-1)$ and $(1,1,-1)$.
Its initial block has length three. Apart from that initial block, its
length-three blocks have indices
\[
 L_j=\floor{\frac{(j+1)(j+2)}4},\qquad j\geq1,
\]
and start at difference indices $2L_j+j$.
At every degree with two absolute maxima the two maximizing exponents
are $b(n)$ and $b(n)+2$, their coefficients have opposite signs,
and $B(n)=M(n)$.
More generally, $B(n)=M(n)$ if and only if either there are two absolute
maxima or $t(n)$ is even.
The signed and absolute modes interlace in the precise sense
$a(n)\in\{b(n)-2,b(n),b(n)+2\}$ for every $n\geq0$.
\end{corollary}

\begin{proof}
Solving $j^2+5j+2\geq2n$ gives
$j\geq\beta_n=(\sqrt{8n+17}-5)/2$. Theorem~\ref{thm:absolute}
therefore gives the formula for $t(n)$ when $n\geq2$; degrees zero and
one give $t=0$ directly. The least residue-class selection gives
$\beta_n\leq b(n)<\beta_n+2$, proving the growth assertion.

For $k\geq1$, the first degree at which $t(n)\geq k$ is
$m_k=2k+h_k$, where
\[
 h_k=\min\{j\geq0:(j+1)(j+2)\geq4k\}.
\]
Indeed, substitute $n=2k+j$ into the threshold in
Theorem~\ref{thm:absolute}. Since $h_1=1$, $m_1=3$.
The increments of $(j+1)(j+2)$ are $2j+4\geq4$, so
$h_{k+1}-h_k\in\{0,1\}$ and $m_{k+1}-m_k\in\{2,3\}$.
On $m_k\leq n<m_{k+1}$, $t(n)=k$; hence the differences are
$m_{k+1}-m_k-1$ copies of $1$ followed by $-1$.
The last index with $h_k=j$ is $L_j$, and every $j\geq1$ occurs,
giving exactly the stated length-three starts.

Strictly decreasing consecutive absolute-coefficient ratios imply that
a tie consists of two consecutive nonzero coefficients. Their signs
alternate in~\eqref{eq:hermite}, so exactly one is positive and attains
the signed maximum. With a unique absolute maximum its sign is
$(-1)^{t(n)}$, proving the last equivalence.
Finally, strict unimodality away from a possible adjacent tie implies
that the largest member with even index $k$ is at $k=t(n)$ when this
index is even, and otherwise at one of the available neighbors
$t(n)-1,t(n)+1$. Their exponents differ from $b(n)$ by two.
\end{proof}

\begin{corollary}\label{cor:values}
For the signed and absolute maximal coefficient sequences A277280 and
A277281 respectively,
\[
 \lim_{n\to\infty}\frac{B(n)}{M(n)}=1,\qquad
 \lim_{n\to\infty}\frac{B(n+1)}{B(n)\sqrt n}
 =\lim_{n\to\infty}\frac{M(n+1)}{M(n)\sqrt n}=\sqrt2.
\]
\end{corollary}
\begin{proof}
If $t=t(n)$ is odd, the ratio of the positive coefficient at exponent
$b(n)+2$ to $M(n)$ is $4t/((b(n)+1)(b(n)+2))$.
Since $t/n\to1/2$ and $b(n)/\sqrt n\to\sqrt2$, this tends to one.
It bounds $B(n)/M(n)$ from below, while one bounds it from above.
For even $t$, $B(n)=M(n)$, so the first limit follows.
Put $j=a(n)$ and $q=q(n)$. Equations~\eqref{eq:closed}
and~\eqref{eq:maxvalue} give, in the two possible transition cases,
\[
 \frac{B(n+1)}{B(n)}=
 \begin{cases}
  2(n+1)/(j+1),&q(n+1)=q,\\
  (n+1)j(j-1)(j-2)/(8(2q+2)(2q+1)),&q(n+1)=q+1.
 \end{cases}
\]
In the second case $j\geq3$. Since $j/\sqrt n\to\sqrt2$
and $q/n\to1/4$, each expression divided by $\sqrt n$ tends to
$\sqrt2$ along its transition case.
Similarly, put $j=b(n)$ and $t=t(n)$. From~\eqref{eq:absolutevalue},
\[
 \frac{M(n+1)}{M(n)}=
 \begin{cases}
  2(n+1)/(j+1),&t(n+1)=t,\\
  (n+1)j/(2(t+1)),&t(n+1)=t+1.
 \end{cases}
\]
Here $j/\sqrt n\to\sqrt2$ and $t/n\to1/2$.
Both cases again give the same limit, completing the proof.
\end{proof}

\section{Coverage and verification}
\begin{center}
\small
\begin{tabular}{@{}>{\raggedright\arraybackslash}p{.42\linewidth}>{\raggedright\arraybackslash}p{.51\linewidth}@{}}
\toprule
Posing item & Mathematical coverage\\
\midrule
A381423, Comments: differences and uniqueness & Proved: Theorem~\ref{thm:signed}; prior difference proof cited above\\
A381423, Comments: block types and growth & Proved: Corollary~\ref{cor:blocks}\\
A381423, Comments: informal spacing & Proved in the exact block-start sense of Corollary~\ref{cor:blocks}; length-four starts need not have constant spacing\\
A277280, definition and ratio plot & Proved: equation~\eqref{eq:maxvalue} and Corollary~\ref{cor:values}\\
A381524, Comments: differences and two maxima & Proved: Theorem~\ref{thm:absolute} and Corollary~\ref{cor:absolute-blocks}\\
A381737, Formula: Conjectures 1--3 & Proved: Theorem~\ref{thm:absolute}, equations~\eqref{eq:gf}--\eqref{eq:recurrence}\\
A277281, definition & Proved: equation~\eqref{eq:absolutevalue} and Corollary~\ref{cor:values}\\
A060821, coefficient formula; A059343, nonzero rows & Prior formulas cited; equation~\eqref{eq:hermite}\\
A060821, auxiliary Riordan, row-sum and related-array identities & Not covered; distinct from the coefficient-mode assertions\\
\bottomrule
\end{tabular}
\end{center}

In the
\href{https://github.com/the-omega-institute/trureturing/blob/a2d277ea0980f1230eff2c7e27838372effccb30/D5/S3/Analytic/Polynomial/HermiteSignedMode.lean}{pinned file \texttt{HermiteSignedMode.lean}},
the declaration \texttt{result}, in the namespace
\texttt{D5.S3.Analytic.Polynomial.}\allowbreak\texttt{HermiteSignedMode},
proves the least-signed-maximizer characterization and the difference
assertion for every natural degree, using the literal expansion~\eqref{eq:hermite}.
Uniqueness, the closed formula, the block and growth formulas, and
the absolute-coefficient results have the written proofs above.
The accompanying \texttt{verify.py} compares these formulas with integer
polynomials generated by $H_{n+1}=2xH_n-2nH_{n-1}$ for $0\leq n\leq600$,
and checks signed and absolute thresholds, transitions, and block indices
through degree one million. It also compares the exact factorial-ratio
identities above through degree 600 and the tie generating function through
coefficient 10,000. These finite comparisons are separate from the all-degree proofs.

\begin{thebibliography}{9}
\small
\bibitem{signed}
M.~Sheppard, OEIS A381423, Comments, revision 8 (6 March 2025;
posed 23 February 2025).
\url{https://oeis.org/history/view?seq=A381423&v=8}.
\bibitem{dlmf}
NIST Digital Library of Mathematical Functions, \S18.5(iii), equation 18.5.13.
\url{https://dlmf.nist.gov/18.5.E13}.
\bibitem{signedvalue}
V.~Reshetnikov, OEIS A277280 (8 October 2016).
\url{https://oeis.org/A277280}.
\bibitem{absolute}
M.~Sheppard, OEIS A381524, Comments (26 February 2025).
\url{https://oeis.org/A381524}.
\bibitem{ties}
M.~Sheppard, OEIS A381737, Formula (5 March 2025).
\url{https://oeis.org/A381737}.
\bibitem{absolutevalue}
V.~Reshetnikov, OEIS A277281 (8 October 2016).
\url{https://oeis.org/A277281}.
\bibitem{triangle}
V.~Jovovic, OEIS A060821, Formula (30 April 2001).
\url{https://oeis.org/A060821}.
\bibitem{nonzero}
N.~J.~A.~Sloane, OEIS A059343 (27 January 2001),
nonzero Hermite coefficient triangle.
\url{https://oeis.org/A059343}.
\end{thebibliography}

\medskip\noindent
\textcopyright\ 2026 Haobo Ma and Wenlin Zhang.
Licensed under Creative Commons Attribution 4.0 International (CC BY 4.0).
\end{document}
